\section{讨论}
\label{sec:comparison-scope}

\paragraph{候选构造的双信号设计。}
关系结构画像描述候选在训练图中的统计接近程度，冻结 B1 分数描述模型已有的关系混淆。\scc{} 将二者作为候选采样依据，并保存逐候选的构造记录，使进入训练的结构与混淆模式可供检查。配对子集中的排序差异与逐题变化提供了进一步分析候选学习的切入点。

\paragraph{训练目标与答案选择。}
候选 InfoNCE 显式训练目标关系与替代关系之间的分数竞争。逐候选评分直接评价这一排序目标，直接解码同时涉及关系文本输出。本文分别报告两种答案选择方式下的表现，以区分给定候选的排序能力与原始解码路径下的预测表现。

\paragraph{固定候选库的适用条件。}
固定候选库将离线构造与 Stage~2 优化分离，使同一运行中的候选来源和采样信号保持固定。其适用条件是一次冻结 B1 全关系评分的预处理成本，以及训练期间固定使用的候选分布。方法面向固定实体对上的关系标签替换，利用训练图统计和关系词表构造排序候选。

\paragraph{比较设计与验证范围。}
主实验对应 NELL23K 的训练种子 2026。B1 与 \scc{} 的同候选重排比较覆盖测试文件前 200 题，采用各自现有训练配置，训练预算存在差异；全量 4,943 题重排结果对应 \scc{}。全量 B1 重排、同预算随机及频率匹配对照、多种子重复与持出正例重叠审计尚未完成。文中的配对差异描述该子集上的模型表现，组合采样的独立作用与跨数据集适用性属于后续验证问题。候选审计按训练已知关系冲突、持出正例重叠和开放世界未判定分别记录，具体规则见附录~\ref{sec:audit-protocol}。
